\section{The Individual Self}

(quoted text from \emph{Indian Philosophy}, Volume 1 by \textbf{Radhakrishnan})


\begin{quotex}
The Upanisads make out that of finite objects the individual self has the highest reality. It comes nearest to the nature of the absolute, though it is not the absolute itself. There are passages where the finite self is looked upon as a reflection of the universe. The whole world is the process of the finite striving to become infinite, and this tension is found in the individual self. According to the Taittiriya the several elements of the cosmos are found in the nature of the individual. In the Chindogya Upanisad (vi. ii. 3 and 4) fire, water and earth are said to constitute the \emph{jivatman} or the individual soul, together with the principle of the infinite.
\end{quotex}


Herein lies the difference between tradition and the modern mind. The latter only recognizes the jivatman, or individual soul, and is unaware of the Infinite which transcends it. So instead of ``striving to become infinite'', it replaces the Infinite by putting its own thoughts in the place of the Infinite.

\begin{quotex}
Man is the meeting-point of the various stages of reality. Prana corresponds to Vayu, the breath of the body to the wind of the world, manas to akasa, the mind of man to the ether of the universe, the gross body to the physical elements. The human soul has affinities with every grade of existence from top to bottom. There is in it the divine element which we call the beatific consciousness, the ananda state, by which at rare moments it enters into immediate relations with the absolute. The finite self or the embodied soul is the Atman coupled with the senses and mind.
\end{quotex}

Here, the terms prana, manas, and the gross body correspond to the Western division of Spirit, soul, and body on the microcosmic level. On the macrocosmic, the Vayu, akasa, and the physical elements are the Holy Spirit, the mind of God or Providence, and the material world.

\begin{quotex}
The different elements are in unstable harmony. ``Two birds, akin and friends, cling to the self-same tree\footnote{\url{https://www.gornahoor.net/?p=1383}}. One of them eats the sweet berry, but the other gazes upon him without eating. In the same tree --- the world tree --- man dwells along with God. With troubles overwhelmed, he faints and grieves at his own helplessness. But when he sees the other, the Lord in whom he delights --- ah, what glory is his, his troubles pass away.'' The natural and the divine have not as yet attained a stable harmony. The being of the individual is a continual becoming, a striving after that which it is not. The infinite in man summons the individual to bring about a unity out of the multiplicity with which he is confronted.
\end{quotex}


In this compact form, we see many themes previously discussed on Gornahoor. For example, the relationship between unity and harmony. There is the disunity in man as he is and lack of harmony of his various parts and thoughts. Man is incomplete; the task of the True Man\footnote{\url{https://www.gornahoor.net/?p=319}} is to realize all his possibilities. The task of the Transcendent man is to realize his unity with the Infinite. Total freedom is achieved only with liberation.

In the next section, we read of the Greater Battle\footnote{\url{https://www.gornahoor.net/?p=3005}}. This battle is fought, first and foremost of all, within a man's own consciousness and involves his whole being. Victory comes only with transcendence.

\begin{quotex}
This tension between the finite and the infinite which is present throughout the world-process comes to a head in the human consciousness. In every aspect of his life, intellectual, emotional and moral, this struggle is felt. He can gain admission into the \textbf{Kingdom of God}, where the eternal verities of absolute love and absolute freedom dwell only by sinking his individuality and transforming the whole of the finite-ness into infiniteness, humanity into divinity. But as finite and human, he cannot reach the fruition or attain the final achievement. The being in which the struggle is witnessed points beyond itself, and so man has to be surpassed. The finite self is not a self-subsisting reality. Be he so, then God becomes only another independent individual, limited by the finite self. The reality of the self is the infinite; the unreality which is to be got rid of is the finite. The finite individual loses whatever reality he possesses if the indwelling spirit is removed. It is the presence of the infinite that confers dignity on the self of man. The individual self derives its being and draws its sustenance from the universal life. \emph{Sub specie aeternitatis}, the self is perfect. There is a psychological side on which the selves repel each other and exclude one another. From this apparent fact of exclusiveness we should not infer real isolation of selves. The exclusiveness is the appearance of distinction. It ought to be referred to the identity, otherwise it becomes a mere abstraction of our minds. The hypothesis of exclusive selves leaves no room for the ideals of truth, goodness and love. These presuppose that man is not perfect as he is, that there is something higher than the actual self which he has to attain to secure peace of mind.
\end{quotex}


Only the man engaged in the greater battle is properly speaking a man. Those not so engaged are still living at the animal level. Victory is the recovery of the Primordial State\footnote{\url{https://www.gornahoor.net/?p=1374}}.

\begin{quotex}
Though the individual soul fighting with the lower nature is the highest in the world, it is not the highest realisable. The striving discordant soul of man should attain to the freedom of spirit, the delight of harmony and the joy of the absolute. Only when the God in him realises itself, only when the ideal reaches its fruition is the destiny of man fulfilled. The struggles, the contradictions and the paradoxes of life are the signs of imperfect evolution, while the harmony, the delight and the peace, mark the perfection of the process of evolution. The individual is the battlefield in which the fight occurs. The battle must be over and the pain of contradiction transcended for the ideal to be realised. The tendency to God which begins in completed man will become then a perfect fruition. Man is higher than all other aspects of the universe, and his destiny is realised when he becomes one with the infinite. Nature has life concealed in it, and when life develops, nature's destiny is fulfilled. Life has consciousness concealed in it, and when it liberates consciousness, its end is reached. The destiny of consciousness is fulfilled when intellect becomes manifest. But the truth of the intellect is reached when it is absorbed in the higher intuition, which is neither thought nor will nor feeling, but yet the goal of thought, the end of will and the perfection of feeling. When the finite self attains the supreme, `the godhead from which it descended, the end of spiritual life is reached. ``When to a man who understands, the self has become all things, what sorrow, what trouble, can there be to him, who has once beheld that unity?''
\end{quotex}



\flrightit{Posted on 2011-10-27 by Cologero}

